\appendix
\chapter{PHỤ LỤC}
\label{app:appendix}

\section{Tái lập thí nghiệm}
\label{app:repro}

Mã nguồn nằm trọn trong gói \texttt{vivlm/} của repo, chạy từ gốc repo theo
thứ tự (chi tiết từng cờ trong \texttt{vivlm/README.md}):

\begin{enumerate}
	\item \texttt{python -m vivlm.tokenizer\_train -{}-parquet <file> -{}-out vivlm/data/tokenizer.json}
	      (huấn luyện BPE; thêm \texttt{-{}-compare} để in bảng fertility so PhoGPT/GPT-2);
	\item \texttt{python -m vivlm.data.prepare\_pretrain -{}-target-tokens 3e9} (tải và
	      token hóa FineWeb2-HQ đến đủ 3 tỷ token);
	\item \texttt{python -m vivlm.pretrain -{}-micro-batch 40} (giai đoạn A; đứt phiên
	      thì thêm \texttt{-{}-resume vivlm/checkpoints/pretrain/latest.pt});
	\item \texttt{python -m vivlm.data.prepare\_sft} rồi
	      \texttt{python -m vivlm.sft -{}-phase projector} và
	      \texttt{python -m vivlm.sft -{}-phase full} (giai đoạn B);
	\item \texttt{python -m vivlm.scst} (giai đoạn C);
	\item \texttt{python -m vivlm.evaluate -{}-mode \{ppl,bpc,caption,vqa\}} và
	      \texttt{python -m vivlm.plot\_logs} (đánh giá, vẽ hình).
\end{enumerate}

Trong khi huấn luyện, \texttt{bash vivlm/watch\_gpu.sh} ghi
\texttt{utilization.gpu} và \texttt{memory.used} mỗi 30 giây vào
\texttt{vivlm/outputs/gpu\_log.csv} --- nguồn của
hình~\ref{fig:gpu-util}. Toàn bộ unit test:
\texttt{python -m pytest vivlm/tests} (52 test, chạy ngoại tuyến trên CPU,
không tải mô hình hay dữ liệu thật).

\section{Danh mục 52 unit test theo nhóm bất biến}
\label{app:tests}

\textbf{Config} (4) --- token đặc biệt đúng thứ tự id 0--3, mặc định kiến
trúc, batch chia hết cho tích lũy gradient; \textbf{GPT} (7) --- RMSNorm
chuẩn hóa đúng, RoPE bảo toàn chuẩn và không xoay vị trí 0, hình dạng
logits/loss, mặt nạ nhân quả (đổi token tương lai không đổi logits quá khứ),
che nhãn $-100$, trói trọng số + đếm đúng 100.682.496 tham số, giải mã tham
lam tất định; \textbf{Tokenizer} (4) --- id token đặc biệt, khứ hồi tiếng
Việt có dấu, token đặc biệt nguyên khối, fertility hợp lệ; \textbf{Dữ liệu
tiền huấn luyện} (5) --- ghi/đọc uint16, mỗi văn bản một \texttt{<|endoftext|>},
cửa sổ lấy mẫu đúng lệch một bước, tái lập theo seed; \textbf{Vòng lặp
pretrain} (4) --- lịch learning rate (đỉnh, đáy, điểm giữa cosine), ghi
checkpoint/log, đúng bộ 7 khóa checkpoint, \emph{khôi phục tương đương từng
bit}; \textbf{Sinh văn bản} (2) --- nạp checkpoint khứ hồi, sinh chuỗi;
\textbf{Hợp nhất} (6) --- pixel-shuffle gom đúng khối $2 \times 2$, hình dạng
projector, loss hữu hạn, che nhãn token ảnh, sinh chỉ trả token mới, chọn
đúng tập tham số theo pha; \textbf{SFT} (2) --- pha projector giảm loss
(quá khớp 4 mẫu), lưu/nạp khứ hồi; \textbf{Dữ liệu SFT} (5) --- che nhãn đúng
vị trí \texttt{<|assistant|>}, cắt ngưỡng, chuẩn hóa ảnh $[-1,1]$, gộp batch
đệm $-100$; \textbf{Đánh giá} (5) --- BLEU/CIDEr khớp hoàn hảo, bất biến
hoa--thường, PPL mô hình ngẫu nhiên, BPC trong khoảng, beam $k{=}1$ trùng
tham lam; \textbf{SCST} (4) --- mặt nạ cắt tại \texttt{<|endoftext|>} đầu
tiên, loss bằng 0 khi lợi thế bằng 0, giá trị loss đúng công thức, gradient
chảy về embedding qua 49 token ảnh; \textbf{Vẽ hình} (3) --- ba loại biểu đồ
xuất file; \textbf{Hồi quy} (1) --- đường dẫn ảnh OpenViVQA giải ra đúng
thư mục gốc (khóa một lỗi thật phát hiện khi rà soát chéo toàn nhánh).

\section{Nhật ký sự cố hạ tầng và cách xử lý}
\label{app:incidents}

\begin{enumerate}
	\item \textbf{\texttt{torch.compile} thất bại vì thiếu trình biên dịch C:}
	      backend Inductor sinh nhân Triton cần \texttt{gcc} tại chỗ; WSL tối
	      giản không có. Xử lý: cài \texttt{gcc}; đo lại cho thấy compile vừa
	      nhanh $2{,}3$ lần vừa tiết kiệm 30\% VRAM so với eager
	      (bảng~\ref{tab:throughput}) --- nếu không cài được, pipeline vẫn chạy
	      với cờ \texttt{-{}-no-compile}.
	\item \textbf{Tràn VRAM âm thầm ở chế độ eager:} WSL cho phép CUDA ``mượn''
	      RAM hệ thống khi hết VRAM --- chương trình không lỗi nhưng chậm đi ba
	      lần. Bài học: theo dõi \emph{mức cấp phát đỉnh} do PyTorch báo chứ
	      không chỉ \texttt{nvidia-smi}.
	\item \textbf{Khôi phục RNG trên GPU:} tensor trạng thái RNG trong
	      checkpoint bị \texttt{map\_location} chuyển lên GPU, trong khi API
	      khôi phục yêu cầu tensor CPU --- lỗi chỉ phát lộ khi resume trên máy
	      GPU thật. Xử lý: ép \texttt{.cpu()} khi khôi phục; khóa bằng test bộ
	      khóa checkpoint.
	\item \textbf{Máy ảo WSL tự tắt theo tiến trình cuối:} mọi phiên dài chạy
	      trong \texttt{tmux} kèm một kết nối giám sát; script khởi chạy kiểm
	      tra checkpoint hiện có để tự nối lại (\texttt{-{}-resume}) sau sự cố.
\end{enumerate}
